\section{离散消息序列信源的统计描述}
\warmth{0}\quad\whisper{序列信源的关键变化：输出不再是一个符号，而是一整段符号向量。}

\begin{strand}
消息序列信源关心的不是单个输出符号，而是一段长度为 \(L\) 的符号序列。
因此统计描述从一维概率分布升级为 \(L\) 维联合概率分布。
\end{strand}

\block{从符号集合到序列集合}
设单个符号的取值集合为 \(\Xcal\)，则长度为 \(L\) 的序列取值集合可写成
\[
\Xcal^L=\Xcal\times \Xcal\times\cdots\times \Xcal,
\]
一次输出的消息序列为
\[
X^L=(X_1,X_2,\cdots,X_L),
\]
它是一个 \(L\) 维随机向量。某个具体序列可写为
\[
x^L=(x_1,x_2,\cdots,x_L).
\]

\begin{definition}[\(L\) 维联合分布]
离散消息序列信源的统计特性由联合概率
\[
\pmf(x^L)=\pmf(x_1,x_2,\cdots,x_L)
\]
描述。一般地，它可以按条件概率链式分解为
\[
\pmf(x_1,x_2,\cdots,x_L)
=\pmf(x_1)\pmf(x_2|x_1)\cdots
\pmf(x_L|x_1,x_2,\cdots,x_{L-1}).
\]
这个式子说明：序列中后一个符号的概率可能依赖于
\answerin[4cm]{前面已经出现的符号}。
\end{definition}

\begin{proposition}[序列取值个数]
若单符号信源有 \(n\) 个可能取值，则长度为 \(L\) 的离散消息序列共有
\[
n^L
\]
个可能序列。
\end{proposition}

\begin{example}
若单符号集合是 \(\{0,1\}\)，长度 \(L=3\)，则全部序列为
\[
000,\ 001,\ 010,\ 011,\ 100,\ 101,\ 110,\ 111.
\]
这里共有 \(2^3=8\) 个可能序列。
\end{example}

\begin{yourturn}
若四进制符号集合为 \(\{0,1,2,3\}\)，序列长度 \(L=2\)，可能序列数为
\[
\answerin[2cm]{16}.
\]
\workspace[2]
\end{yourturn}

\recall{为什么序列信源的概率模型一般要用联合分布，而不是只写每个符号自己的概率？}
